\chapter{Implementation}
\label{ch:implementation}

This chapter documents how the design presented in
Chapter~\ref{ch:analysis-design} was translated into a working system. It
covers the development environment, the two datasets used for the letter and
word recognition tasks, the three classification models that were trained,
and the Android demonstration application that hosts all of them on a single
device. Hyperparameters, library versions and dataset statistics are reported
in tables so that the experiments described in
Chapter~\ref{ch:test-results} are fully reproducible.

%==============================================================================
\section{Development Environment}
\label{sec:impl-environment}

Model development and training were carried out on a single laptop running
Windows~11 with an NVIDIA GeForce RTX~3050 Laptop GPU (4~GB VRAM) and CUDA
toolkit 12.4. The 4~GB VRAM budget was the most binding hardware constraint
of the project: it ruled out larger architectures such as EfficientNet-B3 or
I3D pretrained on Kinetics, and it forced the word model to operate on
precomputed hand landmarks rather than on raw video frames. The same
constraint shaped the deployment target, an off-the-shelf mid-range Android
phone with no dedicated AI accelerator.

The training pipeline is implemented entirely in Python 3.10 inside a virtual
environment. PyTorch~2.6 with the CUDA~12.4 build provides the model and
optimizer code, \texttt{torchvision}~0.21 supplies the EfficientNet-B0
backbone with its ImageNet weights, and Google's MediaPipe~0.10.32 handles
hand-landmark extraction. OpenCV~4.13 is used to decode video files and to
read frames from the webcam during preliminary tests. NumPy and scikit-learn
cover array handling and metric computation, while Matplotlib and Seaborn
generate the confusion matrices and learning curves reported in
Chapter~\ref{ch:test-results}.

The Android application uses a separate but matched stack. The user interface
is written in Kotlin with Jetpack Compose (Material~3), the camera pipeline
is built on CameraX, and hand-landmark inference is delegated to the
MediaPipe Tasks~\textit{HandLandmarker} component in
\textsc{live\_stream} mode. The four classification models are converted to
TensorFlow~Lite and executed through the
\texttt{org.tensorflow:tensorflow-lite} runtime, with the XNNPACK delegate
enabled where applicable. Keeping the training stack in PyTorch and the
deployment stack in TFLite reflects a deliberate split: PyTorch offers the
research conveniences (mixed-precision, eager debugging, rich data tooling),
whereas TFLite is the more mature option for mobile inference on Android.
Table~\ref{tab:impl-environment} summarizes the versions used throughout the
project.

\begin{table}[h]
\centering
\caption{Software stack used for training and deployment.}
\label{tab:impl-environment}
\begin{tabular}{lll}
\toprule
\textbf{Layer} & \textbf{Component} & \textbf{Version} \\
\midrule
Training OS         & Windows                 & 11 (build 26200) \\
GPU runtime         & CUDA                    & 12.4 \\
Training framework  & PyTorch                 & 2.6.0 + cu124 \\
Pretrained models   & torchvision             & 0.21.0 + cu124 \\
Landmark extraction & MediaPipe (Python)      & 0.10.32 \\
Video I/O           & OpenCV                  & 4.13 \\
Numerics            & NumPy                   & 2.4 \\
Metrics             & scikit-learn            & 1.4 \\
\midrule
Mobile language     & Kotlin                  & 1.9 \\
Mobile UI           & Jetpack Compose / M3    & 2025.x BOM \\
Camera              & CameraX                 & 1.3 \\
On-device landmarks & MediaPipe Tasks         & 0.10 \\
On-device inference & TensorFlow Lite (XNNPACK) & 2.16 \\
\bottomrule
\end{tabular}
\end{table}

All training runs were repeatable on a single machine within the time budget
of an academic semester: the longest pipeline, the five-seed BiLSTM ensemble,
trains in under ten minutes once the landmark cache is built. Per-model
training times are reported alongside the corresponding model in
Sections~\ref{sec:impl-efficientnet}--\ref{sec:impl-lstm}.

%==============================================================================
\section{Dataset Preparation}
\label{sec:impl-data}

Two public datasets back the two recognition tasks. The letter task uses the
ASL Alphabet image dataset, which provides clean, frame-level supervision for
the 24 static hand-shapes of American Sign Language. The word task uses a
100-class subset of WLASL, a video-level dataset whose label noise and
collection heterogeneity are well documented in the literature. Both
datasets are described below together with the preprocessing pipelines that
turn them into training-ready tensors.

\subsection{ASL Alphabet (Letter Task)}
\label{sec:impl-data-asl}

The ASL Alphabet dataset hosted on Kaggle contains
\num{87000}~RGB images at 200$\times$200 pixels, equally distributed across
29 classes: the 26 English letters plus three auxiliary classes
(\texttt{del}, \texttt{nothing}, \texttt{space}). Each class therefore
contributes exactly \num{3000} samples, which makes the dataset effectively
balanced. The letters \textit{J} and \textit{Z} are present even though they
involve motion in real ASL; in this dataset they are captured as static
poses, and the system treats them as such.

A stratified 70/15/15 split was applied to the dataset, producing
\num{60900}~training, \num{13050}~validation and \num{13050}~test images
while preserving the original class proportions. During training, mild
augmentation is applied online: horizontal flipping with probability
\num{0.3} to account for left- and right-handed signers, rotations up to
\(\pm 15^{\circ}\), and \texttt{ColorJitter} on brightness and contrast.
Images are resized to 224$\times$224 to match the input resolution expected
by EfficientNet-B0 and normalized with the standard ImageNet statistics.
Validation and test pipelines share the same resize-and-normalize steps but
omit augmentation. For the landmark-based letter model, MediaPipe is applied
once per image and a 63-dimensional vector of joint coordinates is cached
on disk; details of this stage are given in Section~\ref{sec:impl-landmark}.

\subsection{WLASL-100 (Word Task)}
\label{sec:impl-data-wlasl}

The Word-Level American Sign Language (WLASL) dataset~\cite{li2020wlasl}
links each ASL gloss to one or more short YouTube and dictionary-website
clips. The full WLASL2000 release contains roughly \num{21000} videos
across \num{2000} glosses; only the WLASL-100 subset is used in this
project. WLASL-100 keeps the 100 most populous classes and is the canonical
benchmark for compact word recognizers.

Reassembling WLASL-100 in 2026 required more effort than a typical
``download-and-extract'' workflow. Of the eleven source platforms listed in
the original metadata, six were no longer accessible: YouTube videos in
several glosses had been removed, the Flash-based clips from one dictionary
host were unplayable on modern browsers, and one private collection was
unreachable. Two complementary recovery paths were used. First, a
preprocessed Kaggle bundle of WLASL2000 was filtered down to the 100 target
classes; this provided the bulk of clean clips. Second, the
\texttt{url} field in the official WLASL metadata was scraped to retrieve
\num{448} additional videos directly from institutional ASL dictionary
sites (\textsc{aslu}, \textsc{asl5200}, \textsc{valencia-asl},
\textsc{northtexas}), which preserved their content. After deduplication,
\num{1465} videos were retained, a 44\% increase over the
\num{1013}-video baseline that was available at the start of FAZ~6. Per-class
training counts rose from an average of 7.5 to 10.2 videos.

For each video the metadata specifies a sign-active frame window
(\texttt{frame\_start}, \texttt{frame\_end}). The pipeline crops to this
window and then samples 32 frames uniformly along the time axis using
\texttt{np.linspace}; this fixes the temporal resolution regardless of the
original clip length and avoids learning the trivial cue of clip duration.
Each frame is then passed through the MediaPipe \textit{HandLandmarker}
configured with \texttt{num\_hands=2} so that both hands of two-handed signs
can be recovered. The landmark cache is stored once per dataset version and
reused across all training runs.

Two limitations of the recovered dataset are acknowledged here and
revisited in the discussion of Chapter~\ref{ch:test-results}. First, the
remaining sources skew towards dictionary-style content recorded by
professional signers, which under-represents informal everyday signing.
Second, the official WLASL split is person-mixed, so the same signer can
appear in both the training and the test partitions; a stricter
signer-disjoint evaluation is listed as future work in
Chapter~\ref{ch:conclusion}.

\subsection{Landmark Extraction and Feature Encoding}
\label{sec:impl-landmark}

The landmark-based models do not consume raw pixels at training time. Both
the MLP letter classifier and the BiLSTM word classifier work on the
21~3D joint coordinates produced by MediaPipe Hands, which is run once over
the entire dataset as an offline preprocessing pass. Caching landmarks rather
than re-extracting them per epoch keeps GPU utilization high during training
and matches the on-device pipeline, where the mobile application also
receives landmarks from the MediaPipe Tasks component before invoking the
classifier.

For the letter task, MediaPipe is asked for the dominant hand only and the
output is flattened into a single 63-dimensional vector
(\(21\) joints \(\times\) \(3\) coordinates). For the word task the
representation is richer. Both hands are extracted, each is normalized
independently so that the wrist sits at the origin and the distance between
the wrist and the middle-finger MCP joint serves as the unit of length, and
the two hands are routed into fixed ``left'' and ``right'' channels using the
handedness label returned by MediaPipe. A presence flag
\(p \in \{0, 1\}\) is prepended to each hand's channel: if MediaPipe fails
to detect a hand on a particular frame, the channel is zero-filled and its
flag drops to zero. The resulting per-frame feature vector has
\(2 \times (1 + 21 \times 3) = 128\) dimensions, and the per-video tensor has
the shape \((32, 128)\).

Two design choices in this pipeline are worth flagging because they directly
affect the comparison reported in Chapter~\ref{ch:test-results}. First,
using a presence flag instead of a sentinel value lets the LSTM learn an
explicit ``no information here'' signal rather than confusing zero
coordinates with valid positions near the origin. Second, the wrist-relative
normalization removes most of the signer-to-camera distance variance and is
the lightweight stand-in for the body-pose features that a larger
multi-stream model would otherwise rely on.

%==============================================================================
\section{Image-Based Model: EfficientNet-B0}
\label{sec:impl-efficientnet}

The first letter classifier is a standard image-classification pipeline:
EfficientNet-B0~\cite{tan2019efficientnet} pretrained on ImageNet, with its
final classification head replaced by a single fully connected layer that
maps the 1280-dimensional feature embedding to the 29 ASL Alphabet classes.
EfficientNet-B0 was chosen because it offers the best size-versus-accuracy
ratio in its family on commodity hardware: roughly 5.3M parameters,
a 224$\times$224 input resolution that matches the dataset's native aspect
ratio, and a strong ImageNet baseline that transfers well to fine-grained
hand-shape classification.

Training was carried out in two stages, in line with the standard transfer
learning recipe. In the first stage the convolutional backbone was frozen
and only the new classification head was trained for 5~epochs with Adam at
a learning rate of \num{1e-3}. This warm-start step gave the new head time
to align with the pretrained feature space before any backbone weights were
allowed to move. In the second stage all layers were unfrozen and the
network was fine-tuned end-to-end. The fine-tuning stage used the same
optimizer and learning rate; early stopping on validation accuracy halted
training after 6~epochs. The two-stage protocol was a pragmatic choice:
on the available hardware, fine-tuning the backbone immediately from the
ImageNet weights produced unstable gradients during the first epoch, while
the warm-started variant trained smoothly and reached its plateau within a
handful of epochs.

Cross-entropy was used as the loss function, with a batch size of 32 to
fit the 4~GB GPU memory budget alongside the ImageNet-sized activations.
The augmentation pipeline described in
Section~\ref{sec:impl-data-asl} was applied to the training set, and no
augmentation was used on validation or test. End-to-end training took
roughly two hours: 38~minutes for the feature-extraction stage and
85~minutes for the fine-tuning stage. The best checkpoint, selected by
validation accuracy, occupies 15.71~MB when serialized as a PyTorch
\texttt{state\_dict} and runs at \num{490} frames per second on the same
GPU at inference. The same checkpoint is later converted to TensorFlow
Lite for on-device deployment, as described in
Section~\ref{sec:impl-android}.

%==============================================================================
\section{Landmark-Based Letter Model: MLP}
\label{sec:impl-mlp}

The second letter classifier replaces the convolutional backbone with an
explicit feature extractor: MediaPipe Hands detects a single hand in each
image and returns 21 three-dimensional joint coordinates. These 63 numbers
are the only input that the classifier sees; the raw pixels are discarded
entirely. The classifier itself is a four-layer multi-layer perceptron with
hidden widths $[256, 128, 64]$ followed by the 29-way output. Each hidden
layer is followed by batch normalization, a ReLU activation and a dropout
of \num{0.3}. The full network contains roughly \num{60000} parameters and
weighs 0.24~MB once serialized---about 65$\times$ smaller than the
EfficientNet checkpoint.

Landmark extraction is a lossy operation. Of the \num{87000} training
images, MediaPipe was able to recover a hand in
\num{63590} (\SI{73.09}{\percent}); the remaining \num{23410} images were
silently dropped from the landmark-based pipeline. Three sources account
for most of the loss. First, the \texttt{nothing} class contains no hand
by definition and therefore contributes no landmark samples. Second, the
letters~\textit{M} and~\textit{N} involve the thumb tucked under the other
fingers, which routinely confuses the detector on the dataset's
low-resolution crops. Third, motion-blurred or partially out-of-frame
images defeat the detector even when a hand is present. This 73\% recovery
rate is a real cost of the landmark pipeline and is reported alongside the
classifier's accuracy in Chapter~\ref{ch:test-results}; the EfficientNet
model is not affected because it consumes the original pixels regardless of
whether MediaPipe would have succeeded.

Training was a comparatively trivial affair: the network was trained for
100~epochs with Adam at a learning rate of \num{1e-3}, cross-entropy loss,
batch size 64, and no augmentation beyond what is implicit in the landmark
extraction itself. The full training run completed in roughly five minutes
on the same GPU, and the best checkpoint reaches a 99.38\% test accuracy
on the subset of images that survive landmark extraction. The two
classifiers therefore stake out the two ends of the accuracy/footprint
trade-off that is central to this thesis: a 15.71~MB image network that
operates on every image and a 0.24~MB landmark network that depends on a
separate detector for its input. The trade-off is revisited quantitatively
in Chapter~\ref{ch:test-results} and qualitatively in
Chapter~\ref{ch:conclusion}.

%==============================================================================
\section{Word-Level Model: Bidirectional LSTM}
\label{sec:impl-lstm}

Word-level recognition is a fundamentally different problem than letter
recognition. A single frame is insufficient: ASL words encode meaning in
the trajectory of the two hands over roughly a second of video, and a
classifier that ignores this temporal structure cannot do better than
chance on motion-sensitive pairs such as \textit{change}/\textit{how}.
The word model is therefore a recurrent network operating on a
\((32, 128)\) feature tensor produced by the landmark pipeline of
Section~\ref{sec:impl-landmark}: 32 uniformly sampled frames, each frame
holding the two-handed feature vector with its presence flags.

The architecture is a two-layer bidirectional LSTM with a hidden width of
192 and dropout 0.5 between layers, followed by a fully connected
head $(384 \rightarrow 128 \rightarrow 100)$ with a ReLU and a second
dropout step. Bidirectional readout was chosen because ASL words are
short enough that future context within a clip is informative for
disambiguating the early frames; the additional cost is a doubled hidden
state, which the 4~GB GPU absorbs without difficulty. The total parameter
count is roughly \num{938000}, and a serialized checkpoint occupies
5.51~MB. The choice of an LSTM rather than a more recent transformer
encoder was driven by the dataset size: WLASL-100 provides on the order
of ten training videos per class, and the recurrent inductive bias is
better matched to that regime than self-attention is.

\subsection{From V1 to V2: An Iterative Refinement}
\label{sec:impl-lstm-v1v2}

The current model is the second iteration of the word classifier. The
first version, completed in March 2026, used a single-hand 63-dimensional
input and a unidirectional LSTM and reached a validation accuracy of
exactly 20.00\,\%---only marginally better than the 1\,\% random baseline
of a 100-class problem. Diagnosing that failure pointed at three issues
simultaneously, none of which were captured by a single architectural
fix: the input dropped the second hand entirely, video clip boundaries
were not respected, and the model was free to memorize a handful of
clips per class without any regularizer working against it.

V2 addresses each of these issues. The feature representation was widened
to include both hands together with explicit presence flags
(Section~\ref{sec:impl-landmark}); the WLASL \texttt{frame\_start} and
\texttt{frame\_end} fields are now used to crop the sign-active window
before sampling 32~frames; and the training recipe was rewritten around
a regularization-first philosophy. The hidden width was increased from
128 to 192 to give the model enough capacity to exploit the richer input,
and dropout was raised from 0.3 to 0.5. Beyond these architectural
changes, the bulk of the V1$\rightarrow$V2 improvement came from the
training recipe described below.

\subsection{Training Recipe}
\label{sec:impl-lstm-training}

Training is driven by AdamW~\cite{loshchilov2019adamw} with a learning
rate of \num{1e-3} and a weight decay of \num{5e-4}, combined with
\texttt{ReduceLROnPlateau} on validation accuracy
(factor \num{0.5}, patience~5, minimum learning rate \num{1e-5}) and a
gradient-norm clip of \num{1.0}. The loss is cross-entropy with label
smoothing of \num{0.1}~\cite{szegedy2016rethinking}; the smoothing factor
matters here because hand\-edness ambiguity in WLASL makes some training
targets effectively noisy, and a softer target distribution lets the
network commit less aggressively to those examples.

Class imbalance is handled by a \texttt{WeightedRandomSampler} whose
per-sample weight is the inverse of the class frequency in the training
split. The sampler turns every epoch into an effectively balanced pass,
which alone moves validation accuracy from the high 40s into the low
50s. On top of the sampler, training mini-batches are mixed with
Mixup~\cite{zhang2018mixup}: pairs of sequences are linearly interpolated
with $\lambda \sim \mathrm{Beta}(0.4, 0.4)$ and the loss is computed as
the matching convex combination of the two targets. Smaller values of
$\alpha$ produced unstable training curves in early experiments; the
value $0.4$ was selected because it kept training accuracy smooth without
collapsing the class-conditional signal that the model relies on.

The sequence augmentation pack is applied to the training tensors before
Mixup, never to validation or test data. Four operations are layered on
top of each other: a temporal jitter of $\pm 4$ frames with zero
padding, frame-level dropout that fully zeros each frame with probability
\num{0.15}, Gaussian noise of $\sigma = 0.02$ added to the landmark
coordinates (but not to the presence flags), and a time-scale
perturbation that resamples the clip with a uniform stretch factor in
$[0.8, 1.2]$. Together they expose the network to a far broader range of
signing speeds and detector failures than the raw 1465-video dataset
ever could.

Each model is trained for up to 100~epochs with a batch size of 32 and an
early-stopping patience of 15 epochs on validation accuracy. On the
training laptop, a single model converges in under two minutes once the
landmark cache is built, which made the seed-by-seed ablation that
underpins the rest of this section practical to run.

\subsection{Five-Seed Ensemble}
\label{sec:impl-lstm-ensemble}

The final word classifier is an ensemble of five independently trained
LSTM instances. The ensemble shares its architecture, its optimizer, and
its training recipe with the single-model variant; only the random seed
varies, taking the values $\{0, 1, 2, 3, 4\}$. At inference time the
five models' logits are averaged before the softmax, rather than averaging
their softmax outputs:
\[
  \ell_{\text{ens}}(x) = \frac{1}{N} \sum_{n=1}^{N} \ell_n(x),
  \quad
  \hat{y}(x) = \arg\max_c\, \mathrm{softmax}(\ell_{\text{ens}}(x))[c].
\]
Logit averaging preserves the relative magnitudes of each network's
confidence and is the standard choice for small ensembles in the
literature~\cite{dietterich2000ensemble}. The ensemble lifts validation
accuracy from \SI{58.85}{\percent} (seed~0) to \SI{60.08}{\percent}, and
the corresponding test-set Top-1 from \SI{53.23}{\percent} to
\SI{56.22}{\percent} (Chapter~\ref{ch:test-results}). The cost is a
linear blow-up in storage---five 5.51~MB checkpoints together occupy
27.55~MB---and a linear blow-up in inference compute. The mobile demo
exposes the ensemble as a runtime-selectable option alongside the
single-model variant precisely because the size and speed cost may not
be acceptable on every device.

\subsection{Hyperparameter Summary}
\label{sec:impl-lstm-hparams}

Table~\ref{tab:impl-lstm-hparams} consolidates the hyperparameters of the
word model so that they can be reproduced from a single page. Random
seeds, dataset-version pointers and exact training times are deferred to
Section~\ref{sec:impl-repro}.

\begin{table}[h]
\centering
\caption{Hyperparameters of the word recognition model.}
\label{tab:impl-lstm-hparams}
\begin{tabular}{ll}
\toprule
\textbf{Parameter} & \textbf{Value} \\
\midrule
Input shape                & $(32, 128)$ \\
Hidden size                & 192 \\
LSTM layers                & 2, bidirectional \\
Inter-layer dropout        & 0.5 \\
FC head                    & $384 \rightarrow 128 \rightarrow 100$ \\
Trainable parameters       & $\approx \num{938000}$ \\
Optimizer                  & AdamW \\
Learning rate              & \num{1e-3} \\
Weight decay               & \num{5e-4} \\
Loss                       & Cross-entropy, label smoothing 0.1 \\
LR scheduler               & ReduceLROnPlateau ($0.5$, patience 5) \\
Gradient clipping          & $\|\cdot\|_2 \leq 1$ \\
Batch size                 & 32 \\
Max epochs                 & 100 \\
Early-stopping patience    & 15 \\
Temporal jitter            & $\pm 4$ frames \\
Frame dropout              & 0.15 \\
Gaussian noise on coords   & $\sigma = 0.02$ \\
Time-scale range           & $[0.8, 1.2]$ \\
Mixup                      & $\mathrm{Beta}(0.4, 0.4)$ \\
Class balancing            & WeightedRandomSampler ($w_i \propto 1 / n_{y_i}$) \\
Ensemble                   & 5 seeds, logit averaging \\
\bottomrule
\end{tabular}
\end{table}

%==============================================================================
\section{Mobile Demonstration Application}
\label{sec:impl-android}

The four trained models share a single Android application that serves as
the project's deliverable demo and as the platform on which the
comparative claims of Chapter~\ref{ch:test-results} are validated. The
demo lets a user point the device's front camera at one or two hands and
see Top-3 predictions, frame rate and end-to-end latency rendered in
real time. Two classifiers are exposed for letters (EfficientNet-B0 and
the landmark MLP) and two for words (the single-seed BiLSTM and the
five-seed ensemble); each can be selected at runtime through a toggle on
the corresponding screen.

\subsection{Architecture}
\label{sec:impl-android-arch}

The application is written natively in Kotlin with Jetpack Compose for
the user interface, CameraX for the camera pipeline, MediaPipe Tasks for
on-device hand-landmark extraction, and TensorFlow~Lite for classifier
inference. A bottom-bar navigates between three screens: a Home screen
that summarizes the project and the four models, a Letters screen that
hosts the two image/landmark classifiers, and a Words screen that hosts
the two LSTM variants. Dependency wiring is handled by a small
\texttt{AppContainer} singleton instantiated in
\texttt{Application.onCreate}, which keeps the architecture light and
avoids pulling in a full dependency-injection framework for a project of
this size.

Three components do most of the work. \texttt{HandLandmarkerSource}
manages the lifecycle of the MediaPipe \textit{HandLandmarker} in
\textsc{live\_stream} mode and exposes the detected landmarks as a state
flow. \texttt{FrameAnalyzer}, registered as a CameraX
\texttt{ImageAnalysis.Analyzer}, forwards each frame to MediaPipe and
maintains a sliding FPS estimate. \texttt{InferenceManager} owns one
classifier instance per model type and dispatches landmark or pixel
input to the currently selected one. Figure placeholders for the high
level data flow and the per-frame sequence diagram are deferred to
Chapter~\ref{ch:analysis-design}; the discussion here focuses on
implementation choices.

\subsection{Per-Frame Inference Pipeline}
\label{sec:impl-android-pipeline}

CameraX delivers preview frames into an \texttt{ImageAnalysis} queue at
the camera's native rate. Each frame is converted to a bitmap and
submitted to MediaPipe through \texttt{detectAsync}, which returns its
results on a callback thread. The callback runs the active classifier
and pushes the result into a Kotlin flow that the screen's
\texttt{ViewModel} observes. The four classifiers obey a small common
interface but their inputs differ:

\begin{itemize}
  \item \textbf{EfficientNet-B0} consumes the raw bitmap. The frame is
        center-cropped, resized to $224 \times 224$, normalized with
        the ImageNet statistics, and passed to the TFLite runtime,
        which returns a 29-dimensional logit vector. MediaPipe is not
        required for this classifier, but its overlay is still rendered
        so that the two letter models share the same on-screen feedback.
  \item \textbf{Landmark MLP} consumes the 63-dimensional landmark
        vector for the detected hand and returns its own 29-dimensional
        output. If MediaPipe fails to detect a hand in the current
        frame, no inference is performed and the UI displays the last
        confident prediction together with a faded ``no hand visible''
        marker.
  \item \textbf{BiLSTM single} maintains a 32-frame sliding buffer of
        the two-handed 128-dimensional feature vectors described in
        Section~\ref{sec:impl-landmark}. Once the buffer is filled the
        classifier is invoked on every frame, producing a 100-class
        prediction. While the buffer is filling, a sequence-progress
        overlay informs the user that the model is not yet making
        confident predictions.
  \item \textbf{BiLSTM ensemble} runs the five seed checkpoints in
        sequence on the same buffer and averages their logits before
        softmax. The five checkpoints are loaded lazily on the first
        switch to ensemble mode to keep cold-start time predictable.
\end{itemize}

All TFLite interpreters are wrapped in a thin \texttt{TFLiteRunner}
helper that enables the XNNPACK delegate where supported and falls back
to the default CPU kernels otherwise. GPU delegation was experimented
with for the image model on the test device but introduced a long
warm-up period without delivering a stable speedup, so CPU+XNNPACK was
selected as the default.

\subsection{Model Switching and User Interface}
\label{sec:impl-android-ui}

Each camera screen places a segmented toggle in the top-left of the
preview area for switching between its two model variants, an
on-overlay FPS and latency badge in the top-right, and a Top-3
prediction panel anchored at the bottom. Switching models is a
synchronous operation: the toggle updates \texttt{InferenceManager}'s
active classifier and the next analyzer callback uses the new model.
The classifiers are cached across switches so that the second activation
of any model is effectively instantaneous. The Words screen adds a
sequence-progress overlay while the LSTM buffer fills, and the Letters
screen overlays the detected hand-landmark skeleton on top of the
preview as a sanity check that the input the model is consuming is
indeed the user's hand.

Two ergonomic decisions are worth recording. First, the application
defaults to the device's front camera, which is the natural choice for
single-user signing and matches the framing of every training image in
the ASL Alphabet dataset. The landmark overlay is horizontally flipped
to account for the preview's mirror transform so that user-perceived
motion and overlay motion coincide. Second, the preview is rendered
edge-to-edge underneath a transparent status bar, with screen-level
content (titles, toggles) padded into the safe area. The result is a
camera surface that occupies the entire screen rather than the
shrunken-letterbox feel of a standard \texttt{Scaffold} content slot.

\subsection{Edge Cases and Failure Handling}
\label{sec:impl-android-edges}

Three issues recurred during integration and informed the current
defensive behavior. First, the MediaPipe live-stream API requires
monotonically increasing timestamps and silently throws if a stale
frame is submitted after a more recent one has been processed. The
analyzer protects against this with a volatile guard that drops
out-of-order frames before \texttt{detectAsync} is called. Second, on
the test device the MediaPipe GPU delegate took roughly thirty seconds
to warm up while blocking the analyzer thread; the current
implementation pre-loads the \textit{HandLandmarker} on the IO
dispatcher and falls back to the CPU delegate, which initializes in
roughly 300 milliseconds and is sufficient for the full pipeline to
reach above 20~FPS. Third, the application gracefully handles the
permission flow: the first launch displays a permission gate that
re-prompts on rejection and links to system settings, so that a denied
camera permission does not surface as a blank screen or a crash.

%==============================================================================
\section{Reproducibility}
\label{sec:impl-repro}

Every result reported in Chapter~\ref{ch:test-results} can be reproduced
from a clean checkout of the project repository.\footnote{Source code,
training scripts and trained checkpoints are hosted at
\url{https://github.com/yusufkenanakgun/ai-isaret-dili}.} Training the
image classifier requires the ASL Alphabet dataset to be placed in
\texttt{data/asl\_alphabet/}; the landmark cache used by the MLP and the
landmark splits are generated from the same images by
\texttt{src/extract\_kaggle\_landmarks.py}. Training the word classifier
requires the reconstructed WLASL-100 collection described in
Section~\ref{sec:impl-data-wlasl}; the per-clip landmark tensors are
written to \texttt{data/faz6\_v2/landmarks\_wlasl100/} by
\texttt{src/faz6/extract\_v2.py} and consumed by
\texttt{src/faz6/train\_v2.py}.

Random seeds are fixed in each training script: the image and MLP
classifiers use seed~42; the BiLSTM uses seed~0 for the single-model
variant and seeds $\{0, 1, 2, 3, 4\}$ for the ensemble. Library versions
are pinned in \texttt{requirements.txt}; the salient versions are listed
in Table~\ref{tab:impl-environment}. All models were trained on the same
laptop GPU described in Section~\ref{sec:impl-environment}, with the
training times consolidated in Table~\ref{tab:impl-training-summary}.

\begin{table}[h]
\centering
\caption{Training time and final checkpoint size for the four models.}
\label{tab:impl-training-summary}
\begin{tabular}{lcccc}
\toprule
\textbf{Model} & \textbf{Epochs} & \textbf{Training time} & \textbf{Params} & \textbf{Checkpoint} \\
\midrule
EfficientNet-B0          & 5 + 6   & 124 min  & 4.1\,M   & 15.71~MB \\
Landmark MLP             & 100     & 5 min    & 60\,K    & 0.24~MB \\
BiLSTM (single)          & 76--100 & 2 min    & 938\,K   & 5.51~MB \\
BiLSTM ensemble (5 seeds) & 76--100 each & 9.5 min total & 4.7\,M & 27.55~MB \\
\bottomrule
\end{tabular}
\end{table}

The on-device application reads the same checkpoints after they are
converted to TensorFlow~Lite by \texttt{src/convert\_to\_tflite.py}; the
converted files are committed to
\texttt{models/tflite/} for direct consumption by the Android project.
Conversion is deterministic given the PyTorch checkpoints, so the
mobile-side results in Chapter~\ref{ch:test-results} carry the same
reproducibility guarantees as the offline numbers.
